\documentclass[1pt]{article}

\usepackage[a4paper,margin=0.2in]{geometry} % Small margins

\usepackage{multicol} % Multiple columns

\usepackage{amsmath, amssymb} % Math symbols

\usepackage{bm} % Bold math

\usepackage{graphicx} % Insert images

\usepackage{enumitem} % Compact lists

\setlength{\columnsep}{2pt} % Reduce space between columns

\raggedcolumns{}

\begin{document}

\begin{multicols}{1} % 3 columns

\textbf{Buckingham π Theorem}

The Buckingham π theorem states that a physical problem involving $n$ variables and $k$ fundamental dimensions can be reduced to $n-k$ dimensionless parameters $\pi_1, \pi_2, \dots, \pi_{n-k}$. This helps in simplifying complex systems through dimensional analysis.

For a set of variables $x_1, x_2, \dots, x_n$, we express the dimensionless groups as:
\begin{equation*}
    \pi_i = x_1^{a_1} x_2^{a_2} \dots x_n^{a_n}
\end{equation*}
where the exponents $a_i$ are determined by solving a system of linear equations ensuring dimensional consistency.

\textbf{Scaling Relations}

The following relations describe scaling properties:
\begin{equation*}
    \tau \sim d^{1 - \frac{k}{2}}
\end{equation*}

\begin{equation*}
    v \sim \frac{d}{\tau} \sim \frac{d}{d^{1 - \frac{k}{2}}} \sim d^{k/2}
\end{equation*}

\begin{equation*}
    E \sim m v^2 \sim v^2 \sim d^k
\end{equation*}

\begin{equation*}
    L \sim \frac{d^2}{\tau} \sim d^{1 + \frac{k}{2}}
\end{equation*}

\end{multicols}

\end{document}

